\documentclass[addpoints,12pt]{exam}

\usepackage{geometry}
 \geometry{
 a4paper,
 total={170mm,260mm},
 left=20mm,
 top=20mm,
 }

\usepackage{amsmath,amssymb,amsthm}
\DeclareFontEncoding{T1}
\fontencoding{T1}
\usepackage[T2A]{fontenc} 
\usepackage[utf8]{inputenc}
\usepackage[russian]{babel}
\usepackage{graphics,graphicx}
\usepackage{indentfirst}

\usepackage{verbatim}
\usepackage{fancyvrb}  
\usepackage{booktabs}
\usepackage{multirow}

\usepackage{listings}
\usepackage{color}

\renewcommand{\solutiontitle}{\noindent\textbf{Решение:}\par\noindent}

\definecolor{gray}{gray}{0.5}

\lstset{basicstyle=\ttfamily,
    language=Matlab,    
    keepspaces=true,
    extendedchars=\true,
    identifierstyle={ \bf \color{black} },
    showstringspaces=false,
    numbers=left,%
    numbersep=7pt,  
    emph=[1]{case,switch,otherwise,nonlocal,as,yield,zip,print},    
    frame=single,    
    xleftmargin=0.3cm,
    frame=l,framesep=4pt,framerule=0.5pt,
    numberstyle={\scriptsize \color{gray}}
}


\pagestyle{headandfoot}
\runningheadrule
\firstpageheader{ИнтМатПак}{MATLAB 2}{2 марта 2026}
\runningheader{ИнтМатПак}
{MATLAB 2, Стр. \thepage\ из \numpages}
{2 марта 2026}
\firstpagefooter{}{}{}
\runningfooter{}{}{}
\pointname{} 
\boxedpoints
%\pointsinmargin
\pointsinrightmargin


%\header{MATLAB 1}
%}
%{
%	\hfill 29.05.2023
%}

% \footer{Заведующий кафедрой математических методов в экономике}
% {}
% {
%     \hfill проф. Гераськин М. И.
% }


\begin{document}
%\begin{center}
%\fbox{\fbox{\parbox{5.5in}{\centering
%Правильный.}}}
%\end{center}
\vspace{0.1in}
\makebox[\linewidth]{Ф. И. О. \enspace \hrulefill \enspace  группа: \enspace\hrulefill}
\vspace{0.1in}

\begin{questions}

\question[3]
Напишите файл-функцию, которая в качестве аргумента принимает прямоугольную матрицу $A$ размером $n \times m$ случайных чисел, равномерно распределенных в интервале от 0 до 1, и возвращает матрицу размером $n \times 2$. Первый столбец матрицы-результата содержит количество элементов в соответствующей строке матрицы $A$ со значениями меньше 0.5, второй столбец -- количество элементов в строке матрицы $A$ больших 0.9. Напишите два варианта кода файл-функции: с использованием и без использования циклов.
\makeemptybox{7.8cm}

% \question[3] 
% Напишите файл-функцию, которая первым аргументом принимает массив $A$ произвольной длины, заполненный случайными числами, равномерно распределенными в интервале от 0 до 1, и вторым аргументом -- количество отрезков $n \geq 1$, на которые равномерно разбивается интервал [0; 1] (например, для $n=2$ это будут два интервала $[0; 0.5]$ и $(0.5; 1]$). Функция должна возвращать столбец $n \times 1$ с количеством элементов из массива $A$, попадающих в каждый интервал.
% \makeemptybox{8.5cm}

\question[3] 
Некоторая функция $f(t)$ задана таблично на равномерной сетке $t_0, t_1 = t_0+\Delta t, t_2 = t_0+2\Delta t,\ldots,t_n = t_0+n\Delta t$ ($n>1, \Delta t>0$). Напишите функцию, которая в качестве аргументов принимает массив времени $t$, массив значений функции $f$ и возвращает матрицу $n \times 2$, содержащую только те пары  $(t_k, f(t_k))$, для которых выполняется условие убывания: $ f(t_k)-f(t_{k-1})<0 $.

\makeemptybox{9cm}

\question[3] Движение тела по эллиптической орбите в задаче двух тел описывается уравнением Кеплера:
$$
E - e \sin E = M
$$
где $e$ -- эксцентриситет орбиты ($0 \leq e < 1$), $E$ -- эксцентрическая аномалия, $M$ -- средняя аномалия.
Для определения эксцентрической аномалии по известной средней аномалии для орбиты с известным эксцентриситетом может метод простой итерации:
\begin{equation}\label{iter}
    E_{k+1} = e \sin E_k + M, \; k=0,\ldots,n 
\end{equation}
при этом в качестве начального приближения $E_0$ принимают известное значение $M$. 

Напишите файл-функцию, которая принимает в качестве аргументов эксцентриситет $e$, среднюю аномалию $M$ и заданную точность $\varepsilon$ и возвращает эксцентрическую аномалию, вычисленную с заданной точностью. Для вычисления эксцентрической аномалии использовать описанный выше итерационный процесс \eqref{iter} пока $|E_{k+1}-E_{k}|>\varepsilon$.
\makeemptybox{6cm}
\end{questions}

\hqword{Вопрос:}
\hpword{Баллы:}
\hsword{Ответ:}
\htword{Сумма}

\vfill

\begin{center}
\gradetable[h][questions]
\end{center}

\end{document}
